\documentclass[11pt,a4paper]{article}
\usepackage{nesarticle}
\usepackage{tikz}
\hypersetup{pdftitle={Seminar 3: Instructor solutions}}
\begin{document}
\heading{Seminar 3 · Instructor solutions}{Demand, value functions and duality}

\paragraph{A 90-minute route.} Warm-up 15 minutes; problem 1, including
the matrix and general result, 40; problem 2(a)--(f), 30; final questions 5.
The tangential tie in 2(g) is an extension. If time is tight, state repeated
identities after one full derivative calculation per problem.

\section*{Warm-up}
The answers are \textbf{F, T, F, F, F}. Exactly two goods and a common
local point are essential to claims 1 and 3.
\begin{enumerate}[label=\arabic*.]
\item \textbf{False.} By Slutsky,
$x_{i,p_i}=h_{i,p_i}-x_i x_{i,m}$, with $h_{i,p_i}\le0$.
Since $x_i>0$, a Giffen good must be inferior. Both goods cannot be
inferior by the budget argument in claim 3, so they cannot both be Giffen.
\item \textbf{True.} For $u=\sqrt{x_1x_2}$, demand is
$x_i=m/(2p_i)$, so $x_{i,m}=1/(2p_i)>0$ for both goods.
Problem 1 gives another example.
\item \textbf{False.} Local nonsatiation gives budget exhaustion.
Differentiate $p_1x_1+p_2x_2=m$ with respect to income:
\[
 p_1x_{1,m}+p_2x_{2,m}=1.
\]
Both derivatives cannot be negative. If one good is inferior, the other
must be normal. This argument concerns all goods in a two-good budget;
it does not rule out two inferior goods in a larger consumption set.
\item \textbf{False.} At fixed income,
\[
 x_{1,p_1}=\frac{12-m}{(2-p_1)^2}>0,
 \qquad x_{1,m}=-\frac1{2-p_1}<0.
\]
The compensated derivative at the matching target is
\[
 h_{1,p_1}=x_{1,p_1}+x_1x_{1,m}=0.
\]
The law of compensated demand is satisfied; the ordinary demand rises
because of the income effect. This is a genuine utility-maximising demand:
take
\[
 u=\min\{x_1+x_2,\ 2x_1+\tfrac32x_2-6\}.
\]
On the exhausted budget with $p_2=1$, the two scores have slopes
$1-p_1<0$ and $2-3p_1/2>0$ as functions of $x_1$.
Their unique crossing maximises their minimum, and gives the stated demand
and $x_2=(2m-12p_1)/(2-p_1)>0$.
At the crossing $x_1+x_2/2=6$. At its target $\bar u$,
$h=(12-\bar u,2\bar u-12)$ throughout the price region, so the own
compensated derivative really is zero.
\item \textbf{False.} This is a cross-price response, whereas inferiority
is an income response. For $u=\min\{x_1,x_2\}$,
\[
 x_2=\frac m{p_1+p_2},\qquad
 x_{2,p_1}=-\frac m{(p_1+p_2)^2}<0,\qquad
 x_{2,m}=\frac1{p_1+p_2}>0.
\]
Good 2 is normal even though less of it is bought after $p_1$ rises.
\end{enumerate}

\clearpage
\section*{1. Two papers, one research budget}
Here $x_i$ are hours on a paper and $p_i$ are resource costs per hour.
The scores have diminishing marginal returns; they are quality indices,
not probabilities. No additional total-hours constraint is imposed.
For compact notation define
\[
 A(p)=\frac1{p_1}+\frac4{p_2},\qquad D=p_2+4p_1.
\]
\subsection*{(a) Existence, uniqueness and positive quantities}
The positive-price budget is nonempty and compact. Utility is continuous, so
a maximum exists. The weighted sum of square roots is strictly concave on
$\R_+^2$, hence the maximiser on a convex budget is unique. Utility is
strictly increasing in each coordinate, so the whole budget is spent.

Neither axis can contain the optimum. For example, from
$(0,m/p_2)$ buy $\varepsilon>0$ of good 1 and reduce good 2 by
$p_1\varepsilon/p_2$. The gain is $\sqrt{\varepsilon}$ and the loss in the
second term is $O(\varepsilon)$, since its starting quantity is positive.
For small enough $\varepsilon$ the gain dominates. The argument at the other
axis uses $2\sqrt{\varepsilon}$. Do not evaluate ordinary first derivatives
of the square root at zero as if they were finite.

\subsection*{(b) Marshallian demand and indirect utility}
At the interior optimum,
\[
 \frac1{2\sqrt{x_1}}=\lambda p_1,\qquad
 \frac1{\sqrt{x_2}}=\lambda p_2.
\]
Thus $x_1=1/(4\lambda^2p_1^2)$ and
$x_2=1/(\lambda^2p_2^2)$. Exhaustion gives
$A/(4\lambda^2)=m$, so $\lambda=\sqrt A/(2\sqrt m)>0$ and
\[
 \boxed{x_1=\frac{m}{p_1^2A}=\frac{mp_2}{p_1D}},\qquad
 \boxed{x_2=\frac{4m}{p_2^2A}=\frac{4mp_1}{p_2D}}.
\]
Substitution gives
\[
 \boxed{v(p,m)=\sqrt{mA(p)}}.
\]
All quantities are strictly positive for the maintained parameters.

\subsection*{(c) Hicksian demand and expenditure}
Let $z_i=\sqrt{x_i}$. The expenditure problem becomes
\[
 \min_{z_1,z_2\ge0}\ p_1z_1^2+p_2z_2^2
 \quad\text{subject to}\quad z_1+2z_2\ge\bar u.
\]
The objective is strictly convex and coercive and the feasible set is closed
and nonempty. The target binds. With multiplier $\mu$ on
$\bar u-z_1-2z_2\le0$, interior stationarity gives
\[
 2p_1z_1=\mu,\qquad 2p_2z_2=2\mu.
\]
The resulting positive solution is
$\mu=2\bar u/A$, $z_1=\bar u/(p_1A)$ and $z_2=2\bar u/(p_2A)$.
It satisfies the convex programme's sufficient KKT conditions, hence is the
unique minimiser. Squaring,
\[
 \boxed{h_1=\frac{\bar u^2}{p_1^2A^2}
 =\frac{\bar u^2p_2^2}{D^2}},\qquad
 \boxed{h_2=\frac{4\bar u^2}{p_2^2A^2}
 =\frac{4\bar u^2p_1^2}{D^2}}.
\]
Therefore
\[
 \boxed{e(p,\bar u)=\frac{\bar u^2}{A(p)}
 =\frac{\bar u^2p_1p_2}{D}}.
\]
The target is positive, so both quantities and the target multiplier are
positive. The multiplier also equals $e_{\bar u}=2\bar u/A$.

\paragraph{An alternative global proof.} Cauchy--Schwarz gives
\[
 (\sqrt{x_1}+2\sqrt{x_2})^2
 \le (p_1x_1+p_2x_2)\left(\frac1{p_1}+\frac4{p_2}\right).
\]
It bounds utility above by $\sqrt{mA}$ on the budget and bounds the cost of
the target below by $\bar u^2/A$. The displayed choices attain equality.
This proves both values directly without relying on an unverified stationary point.

\subsection*{(d) Roy, Shephard and the value identities}
Write $c_1=1,c_2=4$, so $A_{p_i}=-c_i/p_i^2$. Then
\[
 v_m=\frac{A}{2v}>0,\qquad
 v_{p_i}=-\frac{mc_i}{2vp_i^2},\qquad
 -\frac{v_{p_i}}{v_m}=\frac{mc_i}{p_i^2A}=x_i.
\]
Similarly,
\[
 e_{p_i}=-\frac{\bar u^2 A_{p_i}}{A^2}
 =\frac{\bar u^2c_i}{p_i^2A^2}=h_i.
\]
The duality identities are immediate:
\[
 e(p,v(p,m))=\frac{mA}{A}=m,\qquad
 v(p,e(p,\bar u))=\sqrt{(\bar u^2/A)A}=\bar u.
\]
The positive-target convention removes an absolute-value ambiguity in the
second equality. Substituting $\bar u^2=mA$ into each Hicksian formula gives
$h(p,v(p,m))=x(p,m)$.

\subsection*{(e) Numbers and Slutsky}
At $p=(1,4)$, $m=8$, we have $A=2$, $D=8$,
\[
 x=(4,1),\qquad \bar u=v=4,\qquad e(p,4)=8.
\]
Differentiate Marshallian demand at fixed $m$, and Hicksian demand at fixed
$\bar u$ \emph{before} substituting the matching target:
\[
 x_{1,p_1}=-\frac{mp_2(p_2+8p_1)}{p_1^2D^2}=-6,
 \qquad x_{1,m}=\frac{p_2}{p_1D}=\frac12,
\]
\[
 h_{1,p_1}=-\frac{8\bar u^2p_2^2}{D^3}=-4.
\]
Thus
\[
 \underbrace{-6}_{\text{total}}=
 \underbrace{-4}_{\text{substitution}}
 -\underbrace{4\cdot\tfrac12}_{\text{income correction}}.
\]
Substituting $\bar u=v(p,m)$ into h \emph{before differentiating} would
differentiate ordinary demand and would not compute the compensated derivative.

\subsection*{(f) The Slutsky matrix in this example}
Compute the income correction as an outer product, not componentwise
along the diagonal:
\[
 S=D_px+(D_mx)x^\top
 =\frac{8m}{p_1p_2D^2}
 \begin{pmatrix}-p_2^2&p_1p_2\\p_1p_2&-p_1^2\end{pmatrix}.
\]
Independently, differentiating the Hicksian demands at fixed target gives
\[
 D_ph=\frac{8\bar u^2}{D^3}
 \begin{pmatrix}-p_2^2&p_1p_2\\p_1p_2&-p_1^2\end{pmatrix}.
\]
At $\bar u=v(p,m)$, $\bar u^2=mA=mD/(p_1p_2)$, so the matrices agree.
The displayed matrix is symmetric. Direct multiplication gives $Sp=0$, and
\[
 z^\top Sz=-\frac{8m}{p_1p_2D^2}(p_2z_1-p_1z_2)^2\le0.
\]
At $p=(1,4)$ and $m=8$,
\[
 S=\begin{pmatrix}-4&1\\1&-1/4\end{pmatrix}.
\]
This is negative semidefinite, with a zero price direction. Strict negative
definiteness on the whole space would be the wrong prediction.

\subsection*{(g) From the calculation to the general result}
Under the stated regularity and duality conditions, at a matching target,
\[
 S=D_ph=D^2_{pp}e.
\]
Shephard gives $h_i=e_{p_i}$; Slutsky identifies the derivative matrix with
the corrected Marshallian derivatives. The Cauchy--Clairaut equality of mixed
partials of a $C^2$ function gives symmetry.

To establish the curvature, fix the target-feasible set
$K=\{x\ge0:u(x)\ge\bar u\}$. For positive price vectors p and q and
$0\le t\le1$,
\[
\begin{aligned}
 e(tp+(1-t)q,\bar u)
 &=\min_{x\in K}\{t\,p\cdot x+(1-t)q\cdot x\}\\
 &\ge t\min_{x\in K}p\cdot x+(1-t)\min_{x\in K}q\cdot x.
\end{aligned}
\]
Thus expenditure is concave in prices; convex preferences are not required
for this argument. For any direction z, the restriction
$g(t)=e(p+tz,\bar u)$ is concave near zero, hence
$g''(0)=z^\top D^2_{pp}e\,z\le0$. This proves negative semidefiniteness.

Finally, $h(tp,\bar u)=h(p,\bar u)$: a positive common price multiplier
does not change the cost rankings. Differentiating in t at 1 yields
$D_ph\,p=0$, so $Sp=0$. Summarise the theorem only after the concrete
calculation: it is a property of every regular consumer problem under these
conditions, not a special feature of the square-root utility.

\clearpage
\section*{2. Glitch: coffee, cake, and a head start}
Good 1 is coffee and good 2 is cake. Serving-equivalents are divisible, as
the mathematical domain requires. The head start belongs to the coffee
branch only, including at zero purchases; it is not an additive constant
applied to both branches.
Let
\[
 L(p,m)=1+\frac m{p_1},\qquad R(p,m)=\sqrt{\frac m{p_2}}.
\]
\subsection*{(a) Preferences, the baseline and target geometry}
Utility is continuous. Increasing an active coordinate strictly improves
utility arbitrarily nearby, so preferences are LNS, including at the origin.
It represents complete and transitive preferences. Here $u(0)=1$.

For $\bar u>1$,
\[
 u(x)\ge\bar u\quad\Longleftrightarrow\quad
 x_1\ge\bar u-1\ \ \textbf{or}\ \ x_2\ge\bar u^2.
\]
The set is a union. At target 2, $(1,0)$ and $(0,4)$ reach the target, but
their midpoint $(1/2,2)$ gives $\max\{3/2,\sqrt2\}=3/2<2$.
Preferences are not convex.

\begin{center}
\begin{tikzpicture}[x=2.5cm,y=.65cm]
\fill[NESblue!10] (1,0) rectangle (2,5.5);
\fill[NESblue!10] (0,4) rectangle (1,5.5);
\fill[NESred!10] (0,0)--(1,0)--(0,4)--cycle;
\draw[->] (0,0)--(2.15,0) node[right] {$x_1$};
\draw[->] (0,0)--(0,5.8) node[above] {$x_2$};
\draw[NESblue,thick] (1,0)--(1,4)--(0,4);
\draw[NESred,thick] (1,0)--(0,4);
\fill[NESred] (1,0) circle (1.6pt) node[below right] {$(1,0)$};
\fill[NESred] (0,4) circle (1.6pt) node[above left] {$(0,4)$};
\fill[NESred] (.5,2) circle (1.5pt) node[below left] {midpoint};
\node[font=\small] at (1.45,4.8) {$u\ge2$};
\end{tikzpicture}
\end{center}
The red budget is $4x_1+x_2\le4$; it meets the blue target set only at
the two corners. The example shows why nonunique optima need not include
their midpoint when preferences are nonconvex.

For $\bar u\le1$, zero already attains the target; expenditure is zero and
the unique cost-minimiser is zero. The main calculations use targets above
the baseline. In particular, for $\bar u<1$, the identity $v(p,e)=\bar u$
cannot hold: zero income gives utility 1, not the lower target.

\subsection*{(b) Marshallian demand: compare the two routes}
On the budget,
\[
 x_1+1\le1+m/p_1=L,\qquad \sqrt{x_2}\le\sqrt{m/p_2}=R.
\]
The first bound is attained at $(m/p_1,0)$; the second component attains R
at $(0,m/p_2)$. If $R<1$, that corner's actual utility is 1, and it loses
to $L>1$ anyway. Therefore
\[
 \boxed{v(p,m)=\max\{L(p,m),R(p,m)\}}.
\]
If both quantities are positive, both component bounds are strict, so the
bundle is strictly worse than the largest corner value. There are no other
maximisers. For all positive prices and income,
\[
 \boxed{X(p,m)=
 \begin{cases}
 \{(m/p_1,0)\},&L>R,\\
 \{(m/p_1,0),(0,m/p_2)\},&L=R,\\
 \{(0,m/p_2)\},&L<R.
 \end{cases}}
\]

At $p=(9/2,1)$, let $t=\sqrt m$. Then
\[
 L-R=1+\frac{2t^2}{9}-t
     =\frac29(t-3/2)(t-3).
\]
Thus
\[
\begin{array}{c|c|c}
\text{Income}&\text{optimal route}&v\\\hline
0<m<9/4&\text{good 1}&1+2m/9\\
9/4<m<9&\text{good 2}&\sqrt m\\
m>9&\text{good 1}&1+2m/9
\end{array}
\]
At $m=9/4$, the two corners $(1/2,0)$ and $(0,9/4)$ tie at utility $3/2$.
At $m=9$, $(2,0)$ and $(0,9)$ tie at utility 3. The baseline gives route 1
an initial advantage, the square-root route wins at intermediate incomes,
and the linear route eventually wins again.

\subsection*{(c) Hicksian demand: compare the two costs}
For $\bar u>1$, the cheapest point in the first target branch is
$(\bar u-1,0)$ at cost $C_1=p_1(\bar u-1)$; in the second it is
$(0,\bar u^2)$ at cost $C_2=p_2\bar u^2$. Hence
\[
 \boxed{e(p,\bar u)=\min\{C_1,C_2\}},\qquad
 \boxed{H(p,\bar u)=
 \begin{cases}
 \{(\bar u-1,0)\},&C_1<C_2,\\
 \{(\bar u-1,0),(0,\bar u^2)\},&C_1=C_2,\\
 \{(0,\bar u^2)\},&C_1>C_2.
 \end{cases}}
\]
Each branch has only its corner as the cheapest point at positive prices,
so the tie adds exactly two bundles, not their convex combinations.

At $p=(9/2,1)$,
\[
 C_1-C_2=\frac92(\bar u-1)-\bar u^2
        =-(\bar u-3/2)(\bar u-3).
\]
For $1<\bar u<3/2$ and for $\bar u>3$, choose good 1. For
$3/2<\bar u<3$, choose good 2. At each equality choose either corner.
The expenditure values at the two switching targets are $9/4$ and 9.
For all targets $\bar u\le1$, $e=0$ and $H=\{(0,0)\}$.

\subsection*{(d) Roy and Shephard within strict regimes}
In a strict first-route regime, $v=1+m/p_1$ and
\[
 v_m=1/p_1,\qquad v_{p_1}=-m/p_1^2,\qquad v_{p_2}=0,
\]
so Roy gives $(m/p_1,0)$. In a strict second-route regime, $v=\sqrt{m/p_2}$,
\[
 v_m=\frac1{2\sqrt{mp_2}},\quad
 v_{p_1}=0,\quad v_{p_2}=-\frac{\sqrt m}{2p_2^{3/2}},
\]
and Roy gives $(0,m/p_2)$.

If $C_1<C_2$, $e=p_1(\bar u-1)$ and
$\nabla_pe=(\bar u-1,0)$. If $C_2<C_1$, $e=p_2\bar u^2$ and
$\nabla_pe=(0,\bar u^2)$. Shephard recovers the corresponding Hicksian
corner. At a tie the price gradients of the two value branches differ,
so the full gradient needed by Roy or Shephard is unavailable.

\subsection*{(e) Duality and matching the switches}
Set $U=v(p,m)$. Every affordable bundle has
$x_1\le m/p_1\le U-1$ and $x_2\le m/p_2\le U^2$.
More directly, $U\ge L$ implies $p_1(U-1)\ge m$, and $U\ge R$ implies
$p_2U^2\ge m$. At least one equality holds, hence
\[
 e(p,v(p,m))=\min\{p_1(U-1),p_2U^2\}=m.
\]
The equality branches are exactly the winning Marshallian routes, so the
demand sets match, including at ties.

Conversely, set $M=e(p,\bar u)$. Since
$M\le p_1(\bar u-1)$ and $M\le p_2\bar u^2$,
\[
 1+M/p_1\le\bar u,\qquad \sqrt{M/p_2}\le\bar u.
\]
At least one equality holds, so $v(p,M)=\bar u$. The corners attaining
that equality are exactly $H(p,\bar u)$.

At the specified prices the boundaries pair as
\[
 m=9/4\ \longleftrightarrow\ \bar u=3/2,\qquad
 m=9\ \longleftrightarrow\ \bar u=3.
\]
No convexity of preferences is required for these duality statements.

\subsection*{(f) The two exact boundary examples}
At $p=(9/2,1)$, $m=9/4$ and $\bar u=3/2$,
\[
 X(p,9/4)=H(p,3/2)=\{(1/2,0),(0,9/4)\}.
\]
The midpoint $(1/4,9/8)$ gives
$\max\{5/4,3/(2\sqrt2)\}=5/4<3/2$ and fails the target.
At the other boundary,
\[
 X(p,9)=H(p,3)=\{(2,0),(0,9)\}.
\]
The midpoint $(1,9/2)$ gives $\max\{2,3/\sqrt2\}=3/\sqrt2<3$.
Nonuniqueness persists in the corresponding expenditure problem, and the
midpoint is not optimal in either problem.

\subsection*{(g) Extension: a tangential tie}
At fixed prices $(4,1)$,
\[
 1+\frac m4-\sqrt m=\frac14(\sqrt m-2)^2\ge0.
\]
Thus $v((4,1),m)=1+m/4$ for every $m>0$, and $v_m=1/4$ even at $m=4$.
At that income, however, $(1,0)$ and $(0,4)$ are both optimal.

There is no contradiction with Roy: differentiability in income alone is
insufficient. At the tie the price gradients are $(-1/4,0)$ on the first
branch and $(0,-1)$ on the second. The full indirect utility is not
differentiable in prices there. Expenditure at fixed prices is likewise
smooth in the target at 2, but its two price gradients are $(1,0)$ and $(0,4)$.
Do not claim that every nonunique optimum forces a kink in the income-only
value graph: this tie is a counterexample.

\paragraph{General switch geometry, for the teaching assistant.}
Writing $r=p_1/p_2$, the tie equation in utility is
$\bar u^2-r\bar u+r=0$. If $r<4$ the first route always wins. At $r=4$
there is one tangential tie at $\bar u=2$, corresponding to $m=4p_2$.
If $r>4$, there are two thresholds
\[
 \bar u_{\pm}=\frac{r\pm\sqrt{r^2-4r}}2>1,\qquad
 m_{\pm}=p_2\bar u_{\pm}^2=p_1(\bar u_{\pm}-1).
\]
The second route wins strictly between the thresholds and the first outside.
The comparison formulas in (b)--(c) cover all prices without requiring
students to carry this discriminant calculation throughout the main exercise.
\end{document}
